% GENERATED FILE. Edit canonical lessons, then run npm run generate:books.
% canonical-source-hash: fnv1a64:718c27a3f50be2ab
% canonical-lessons: LA-C442-salum, LA-C442-pelagus, LA-C442-aequor, LA-C442-marinus, LA-C442-maritimus

\chapter{Open sea, High sea, Level sea, Of the sea, Coastal}
\label{ch:la-open-sea-high-sea-level-sea-of-the-sea-coastal}
\hlchaptermodality{442}

\begin{chapteropening}
\textbf{By the end of this chapter:} \emph{I can talk about the open sea and the coast: the open sea, the high sea, the level sea, things of the sea, and coastal places.}

Complete the last lesson of chapter 442: I can talk about the open sea and the coast: the open sea, the high sea, the level sea, things of the sea, and coastal places.
\end{chapteropening}

\section[salum, salī]{salum, salī — the open sea, the swell (of the sea)}

\label{lesson:LA-C442-salum}



\begin{quote}
Before the new one: say the Latin for a lighthouse, then the Latin for a headland.
\end{quote}

\subsection*{You'll want to know: salum, salī}
\textbf{salum, salī} — \textquotedblleft{}the open sea, the swell (of the sea)\textquotedblright{}.

\begin{quote}
Nāvēs extrā portum in salō ad ancorās stant.
\end{quote}

\begin{quote}
\emph{The ships lie at anchor on the open sea outside the harbour.}
\end{quote}

Second declension, neuter: \emph{salum, salī n.} The rolling sea off the coast, where ships ride at anchor.

\begin{cousinweb}
\emph{salum} matches Greek \emph{salos} 'the rolling swell of the sea' (Lewis \& Short), but EtymDB gives no parent, so whether it was borrowed is not settled.
\end{cousinweb}

\begin{grammarlens}[title={What you've built}]
One more word for telling, explaining and describing at B1.
\end{grammarlens}

\subsection*{Your turn}
\begin{itemize}
\raggedright
  \item \emph{Say it:} \emph{salum, salī}
  \item \emph{Say it:} \emph{salum, salī}, once more
  \item \emph{Say it:} a sentence of your own with \emph{salum, salī}
  \item \emph{Recall:} say the Latin for a lighthouse, then the Latin for a headland, then say \emph{salum, salī} again
\end{itemize}

\begin{tcolorbox}[breakable,colback=teal!4,colframe=teal!35!black,title={Before you move on}]
What does salum, salī mean? (\textquotedblleft{}The open sea, the swell (of the sea)\textquotedblright{}.) Say it once more.
\end{tcolorbox}
\section[pelagus, pelagī]{pelagus, pelagī — the high sea, the open sea}

\label{lesson:LA-C442-pelagus}



\begin{quote}
Before the new one: say the Latin for a headland, then the Latin for the open sea.
\end{quote}

\subsection*{You'll want to know: pelagus, pelagī}
\textbf{pelagus, pelagī} — \textquotedblleft{}the high sea, the open sea\textquotedblright{}.

\begin{quote}
Nāvis parva in pelagō sōla est.
\end{quote}

\begin{quote}
\emph{The small ship is alone on the high sea.}
\end{quote}

Second declension but neuter: \emph{pelagus, pelagī n.} A poetic and Greek word. English \emph{pelagic} 'of the open sea' is related.

\begin{cousinweb}
\emph{pelagus} was borrowed from Greek \emph{pelagos} 'the sea, the open sea' (EtymDB; Lewis \& Short), which explains why it keeps the Greek neuter gender in \emph{-us}.
\end{cousinweb}

\begin{grammarlens}[title={What you've built}]
One more word for telling, explaining and describing at B1.
\end{grammarlens}

\subsection*{Your turn}
\begin{itemize}
\raggedright
  \item \emph{Say it:} \emph{pelagus, pelagī}
  \item \emph{Say it:} \emph{pelagus, pelagī}, once more
  \item \emph{Say it:} a sentence of your own with \emph{pelagus, pelagī}
  \item \emph{Recall:} say the Latin for a headland, then the Latin for the open sea, then say \emph{pelagus, pelagī} again
\end{itemize}

\begin{tcolorbox}[breakable,colback=teal!4,colframe=teal!35!black,title={Before you move on}]
What does pelagus, pelagī mean? (\textquotedblleft{}The high sea, the open sea\textquotedblright{}.) Say it once more.
\end{tcolorbox}
\section[aequor, aequoris]{aequor, aequoris — the level sea, the sea's surface; a level plain}

\label{lesson:LA-C442-aequor}



\begin{quote}
Before the new one: say the Latin for the open sea, then the Latin for the high sea.
\end{quote}

\subsection*{You'll want to know: aequor, aequoris}
\textbf{aequor, aequoris} — \textquotedblleft{}the level sea, the sea's surface; a level plain\textquotedblright{}.

\begin{quote}
Ventus nūllus est, et aequor quiētum est.
\end{quote}

\begin{quote}
\emph{There is no wind, and the surface of the sea is still.}
\end{quote}

Third declension, neuter: \emph{aequor, aequoris n.} First any flat surface, then above all the calm, level sea (a word of poetry).

\begin{cousinweb}
\emph{aequor} is built on \emph{aequus} 'level, even' (EtymDB; Lewis \& Short): a level surface, and so the sea when it is calm.
\end{cousinweb}

\begin{grammarlens}[title={What you've built}]
One more word for telling, explaining and describing at B1.
\end{grammarlens}

\subsection*{Your turn}
\begin{itemize}
\raggedright
  \item \emph{Say it:} \emph{aequor, aequoris}
  \item \emph{Say it:} \emph{aequor, aequoris}, once more
  \item \emph{Say it:} a sentence of your own with \emph{aequor, aequoris}
  \item \emph{Recall:} say the Latin for the open sea, then the Latin for the high sea, then say \emph{aequor, aequoris} again
\end{itemize}

\begin{tcolorbox}[breakable,colback=teal!4,colframe=teal!35!black,title={Before you move on}]
What does aequor, aequoris mean? (\textquotedblleft{}The level sea, the sea's surface; a level plain\textquotedblright{}.) Say it once more.
\end{tcolorbox}
\section[marīnus, marīna, marīnum]{marīnus, marīna, marīnum — of the sea, marine}

\label{lesson:LA-C442-marinus}



\begin{quote}
Before the new one: say the Latin for the high sea, then the Latin for the level sea.
\end{quote}

\subsection*{You'll want to know: marīnus, marīna, marīnum}
\textbf{marīnus, marīna, marīnum} — \textquotedblleft{}of the sea, marine\textquotedblright{}.

\begin{quote}
Aqua marīna salsa est, nōn dulcis.
\end{quote}

\begin{quote}
\emph{Sea water is salty, not sweet.}
\end{quote}

An adjective of three endings, from \emph{mare} 'sea'. English \emph{marine} and \emph{marina} come from it.

\begin{cousinweb}
\emph{marīnus} is built on \emph{mare} 'sea', inherited from Proto-Italic \emph{mari}, PIE \emph{móri} (EtymDB; Lewis \& Short). English \emph{marine} and \emph{marina} come from it.
\end{cousinweb}

\begin{grammarlens}[title={What you've built}]
One more word for telling, explaining and describing at B1.
\end{grammarlens}

\subsection*{Your turn}
\begin{itemize}
\raggedright
  \item \emph{Say it:} \emph{marīnus, marīna, marīnum}
  \item \emph{Say it:} \emph{marīnus, marīna, marīnum}, once more
  \item \emph{Say it:} a sentence of your own with \emph{marīnus, marīna, marīnum}
  \item \emph{Recall:} say the Latin for the high sea, then the Latin for the level sea, then say \emph{marīnus, marīna, marīnum} again
\end{itemize}

\begin{tcolorbox}[breakable,colback=teal!4,colframe=teal!35!black,title={Before you move on}]
What does marīnus, marīna, marīnum mean? (\textquotedblleft{}Of the sea, marine\textquotedblright{}.) Say it once more.
\end{tcolorbox}
\section[maritimus, maritima, maritimum]{maritimus, maritima, maritimum — coastal, by the sea; of the sea}

\label{lesson:LA-C442-maritimus}



\begin{quote}
Before the new one: say the Latin for the level sea, then the Latin for of the sea.
\end{quote}

\subsection*{You'll want to know: maritimus, maritima, maritimum}
\textbf{maritimus, maritima, maritimum} — \textquotedblleft{}coastal, by the sea; of the sea\textquotedblright{}.

\begin{quote}
Oppida maritima multōs mercātōrēs habent.
\end{quote}

\begin{quote}
\emph{Towns by the sea have many merchants.}
\end{quote}

An adjective of three endings, from \emph{mare} 'sea'. \emph{Ōra maritima}, the sea coast. English \emph{maritime} comes from it.

\begin{cousinweb}
\emph{maritimus} is built on \emph{mare} 'sea' with an old suffix also seen in \emph{fīnitimus} 'neighbouring' (EtymDB; Lewis \& Short). English \emph{maritime} comes from it.
\end{cousinweb}

\begin{grammarlens}[title={What you've built}]
One more word for telling, explaining and describing at B1.
\end{grammarlens}

\subsection*{Your turn}
\begin{itemize}
\raggedright
  \item \emph{Say it:} \emph{maritimus, maritima, maritimum}
  \item \emph{Say it:} \emph{maritimus, maritima, maritimum}, once more
  \item \emph{Say it:} a sentence of your own with \emph{maritimus, maritima, maritimum}
  \item \emph{Recall:} say the Latin for the level sea, then the Latin for of the sea, then say \emph{maritimus, maritima, maritimum} again
\end{itemize}

\begin{tcolorbox}[breakable,colback=teal!4,colframe=teal!35!black,title={Before you move on}]
What does maritimus, maritima, maritimum mean? (\textquotedblleft{}Coastal, by the sea; of the sea\textquotedblright{}.) Say it once more.
\end{tcolorbox}
